\documentclass[11pt]{article}
\usepackage[a4paper,margin=27mm]{geometry}
\usepackage{amsmath,amssymb,amsthm,booktabs,array}
\usepackage[T1]{fontenc}
\usepackage{lmodern}
\usepackage[colorlinks=true,linkcolor=blue,citecolor=blue,urlcolor=blue]{hyperref}
\newtheorem{theorem}{Theorem}
\newtheorem{lemma}[theorem]{Lemma}
\newtheorem{corollary}[theorem]{Corollary}
\newcommand{\coeff}[2]{[x^{#1}]\,#2}
\newcommand{\Q}{\mathbb{Q}}
\newcommand{\Z}{\mathbb{Z}}
\title{Proof of Hanna's square-exponent factorial integrality conjectures}
\author{Haobo Ma\textsuperscript{1}\qquad Wenlin Zhang\textsuperscript{2,1}\\[3pt]
  \small\textsuperscript{1}ChronoAI Pte Ltd, Singapore\\
  \small\textsuperscript{2}National University of Singapore}
\date{November 17, 2026}
\begin{document}
\maketitle
\begin{abstract}
In the COMMENT sections of OEIS A350410 and A350411 (revision 5) and
A350412 (revision 9), Paul D. Hanna conjectured that the three sequences
consist entirely of integers; we prove all three conjectures.
Their factorial denominators cancel by a coefficient-divisibility identity.
For each series we also establish existence and uniqueness over the rationals.
More generally, a degree-dependent cutoff $M_n\geq1$ gives the same conclusions
whenever $\gcd(n,M_n^2)$ divides $M_n$, with any prescribed integer linear
coefficient. This arithmetic condition is necessary and sufficient for
the last denominator to cancel for every integer polynomial, and the
solutions depend homogeneously on the prescribed linear coefficient.
All defining sums are finite sums of formal power series.
\end{abstract}

\section{The conjectures and a common theorem}
Hanna's three entries \cite{hanna410,hanna411,hanna412} each state:
``This sequence is conjectured to consist entirely of integers.'' We prove all three
assertions. Their defining equations determine a new coefficient using a finite
factorial-weighted sum of square powers; the issue is why that rational
recursion always produces an integer. The decisive cancellation is a
divisibility property of the last summand.

For a formal power series $F\in\Q[[x]]$, write $[x^n]F$ for its coefficient of
$x^n$. For an integer $M\geq1$, define the finite sum
\[
 S_M(F)=\sum_{m=0}^{M}\frac{F^{m^2}}{m!}.
\]
The term with $m=0$ is $1$, including when $F=0$.
The examples in A350410 and A350411 specify the initial terms $-1+x$.
Their NAME equations are, respectively,
\begin{equation}\label{eq:two}
 [x^n]S_n(B)=0,\qquad [x^n]S_{n+1}(B)=0\qquad(n\geq2).
\end{equation}
The series $A$ in A350412 starts with $x$ and satisfies
\begin{equation}\label{eq:third}
 A(0)=0,\quad [x]A=1,\quad [x^n]S_{2n}(A-1)=0\qquad(n\geq2).
\end{equation}
Thus $B=A-1$ has the same normalization as the first two series.
The prescribed linear coefficient matters: these equations impose no condition
in degree one.

\begin{theorem}\label{thm:main}
Let $(M_n)_{n\geq2}$ be any sequence of positive integers, and let $t\in\Q$.
There is a unique series
\[
 B=-1+tx+\sum_{n\geq2}b_nx^n\in\Q[[x]]
 \quad\text{such that}\quad [x^n]S_{M_n}(B)=0\quad(n\geq2).
\]
If $t\in\Z$ and
\begin{equation}\label{eq:criterion}
 \gcd(n,M_n^2)\mid M_n\qquad(n\geq2),
\end{equation}
then $B\in\Z[[x]]$.
\end{theorem}

\section{Cancellation of the factorial denominator}
We first compute the coefficient multiplying the unknown $b_n$.
If $P$ has constant term $-1$ and degree less than $n$, then for every
integer $r\geq1$ and every $u\in\Q$,
\[
 [x^n](P+ux^n)^r=[x^n]P^r+r(-1)^{r-1}u.
\]
Indeed, terms containing two factors $ux^n$ have degree at least $2n>n$,
and the remaining $r-1$ factors in the linear term contribute their constants.
Consequently
\begin{equation}\label{eq:slope}
 [x^n]S_M(P+ux^n)=[x^n]S_M(P)+\lambda_M u,
 \qquad
 \lambda_M=\sum_{m=1}^{M}\frac{(-1)^{m-1}m^2}{m!}
          =\frac{(-1)^{M-1}}{(M-1)!}.
\end{equation}
Here $m^2$ and $m$ have the same parity. To obtain the last equality, put
$h_0=0$ and $h_m=(-1)^{m-1}/(m-1)!$ for $m\geq1$.
For $m=1$ the summand is $h_1-h_0=1$; for $m\geq2$ it is
\[
 \frac{(-1)^{m-1}m}{(m-1)!}
 =\frac{(-1)^{m-1}}{(m-1)!}
  -\frac{(-1)^{m-2}}{(m-2)!}=h_m-h_{m-1}.
\]
The sum therefore telescopes. In particular, $\lambda_M\ne0$.

\begin{lemma}\label{lem:div}
For $P\in\Z[[x]]$, $n\geq1$ and $r\geq1$,
\[
 \frac{r}{\gcd(n,r)}\ \bigm|\ [x^n]P^r.
\]
In particular, if $M\geq1$ and $\gcd(n,M^2)\mid M$, then
$M\mid[x^n]P^{M^2}$.
\end{lemma}
For nonnegative coefficients, the divisibility statement is Eger's
Theorem~19 \cite{eger}. Reducing coefficients modulo $r/\gcd(n,r)$ gives
the signed version; the following direct proof covers both.
\begin{proof}
Formal differentiation gives
\[
 n[x^n]P^r=[x^{n-1}](P^r)'=r[x^{n-1}]P'P^{r-1}.
\]
Thus $r$ divides $n[x^n]P^r$. Dividing $n$ and $r$ by their greatest common
divisor and applying coprimality gives the first assertion.
For the second, set $r=M^2$ and $d=\gcd(n,M^2)$. If $d\mid M$, then
$M^2/d=M(M/d)$ is a multiple of $M$, so the first assertion applies.
\end{proof}

\begin{proof}[Proof of Theorem~\ref{thm:main}]
Set $b_0=-1$ and $b_1=t$. Suppose $b_0,\ldots,b_{n-1}$ are known, where
$n\geq2$, and put $P_n=\sum_{j=0}^{n-1}b_jx^j$ and $M=M_n$.
Equation~\eqref{eq:slope} forces precisely the value
\begin{equation}\label{eq:recursion}
 b_n=(-1)^M(M-1)!\,[x^n]S_M(P_n).
\end{equation}
This defines one rational coefficient at each step. Higher coefficients cannot
affect a coefficient of degree $n$, so the resulting formal series satisfies
every required equation. The same formula forces the coefficients of any other
solution one by one, proving uniqueness.

Now assume $t\in\Z$ and \eqref{eq:criterion}, and argue by induction that
$P_n\in\Z[x]$. Multiplying the right-hand coefficient by $(M-1)!$ gives
\begin{equation}\label{eq:clear}
 (M-1)!\,[x^n]S_M(P_n)
 =\sum_{m=0}^{M-1}\frac{(M-1)!}{m!}[x^n]P_n^{m^2}
   +\frac{[x^n]P_n^{M^2}}{M}.
\end{equation}
Every factorial ratio in the sum is an integer, and the final quotient is an
integer by Lemma~\ref{lem:div}. Formula~\eqref{eq:recursion} therefore gives
$b_n\in\Z$, completing the induction. This also covers $M=1$: the sum in
\eqref{eq:clear} has only its $m=0$ term, and the final divisor is $1$.
\end{proof}

\begin{corollary}\label{cor:hanna}
Each of the three normalized series in
\eqref{eq:two}--\eqref{eq:third} exists, is unique over $\Q$, and has only
integer coefficients.
More generally, the integrality conclusion holds if, independently in each
degree $n$, the cutoff $M_n$ is any positive multiple of $n$, or is coprime
to $n$, and the linear coefficient is any integer.
\end{corollary}
\begin{proof}
If $M=kn$ with $k\geq1$, then $\gcd(n,M^2)=n\mid M$.
If $\gcd(n,M)=1$, then $\gcd(n,M^2)=1\mid M$.
These facts apply to $M_n=n$, $n+1$ and $2n$.
Take $t=1$ in Theorem~\ref{thm:main}; for the third series translate $B$ to
$A=B+1$, which preserves all positive-degree coefficients and changes the
constant term to zero.
\end{proof}

\section{Scaling and sharpness of the cancellation condition}
\begin{corollary}\label{cor:scaling}
Fix any positive cutoff sequence $(M_n)_{n\geq2}$, and let $B_t$ be the
unique solution in Theorem~\ref{thm:main} with linear coefficient $t\in\Q$.
Then $B_t(x)=B_1(tx)$. In particular,
$[x^n]B_t=t^n[x^n]B_1$ for every $n\geq0$.
If the gcd condition holds, every coefficient is an integer multiple of
$t^n$ whenever $t\in\Z$.
\end{corollary}
\begin{proof}
Substitution $x\mapsto tx$ multiplies each degree-$n$ coefficient of
$S_{M_n}(B_1)$ by $t^n$, preserving its vanishing for $n\geq2$.
It gives the required constant and linear coefficients, so uniqueness
identifies it with $B_t$. This also applies to $t=0$.
\end{proof}

\begin{theorem}\label{thm:sharp}
For positive integers $n,M$, the following are equivalent:
\begin{enumerate}
\item $\gcd(n,M^2)\mid M$;
\item for every $P\in\Z[x]$ with $P(0)=-1$, one has
$M\mid[x^n]P^{M^2}$.
\end{enumerate}
Equivalently, for every prime $p\mid M$, the exponent of $p$ in $n$
is at most its exponent in $M$.
\end{theorem}
\begin{proof}
Lemma~\ref{lem:div} proves the forward implication.
Write $v_p(a)$ for the exponent of a prime $p$ in a positive integer $a$,
extended to nonzero rationals by subtraction of numerator and denominator
exponents. If (1) fails, some $p\mid M$ has
$e=v_p(M)\geq1$ and $v_p(n)\geq e+1$.
Set $N=M^2$, $k=p^{e+1}$, $s=n/k$ and $P=-1+x^s$.
Then $s$ is a positive integer and $k\leq p^{2e}\leq N$, while
\[
 [x^n]P^N=(-1)^{N-k}\binom Nk,
 \qquad \binom Nk=\frac Nk\binom{N-1}{k-1}.
\]
For $1\leq i<k$, one has $v_p(i)\leq e<2e=v_p(N)$, whence
$v_p(N-i)=v_p(i)$. The product
$\binom{N-1}{k-1}=\prod_{i=1}^{k-1}(N-i)/i$ therefore has valuation zero.
It follows that $v_p\binom Nk=2e-(e+1)=e-1$; hence $M$ does not divide
the displayed coefficient. This proves the converse.
Finally, $v_p(\gcd(n,M^2))=\min(v_p(n),2v_p(M))$ gives the stated
prime-by-prime equivalence.
\end{proof}

Theorem~\ref{thm:sharp} concerns cancellation for every integer polynomial;
it does not assert necessity for a particular recursively generated series.
The unrestricted positive-cutoff extension of integrality is not asserted
in Hanna's entries and fails: for the constant cutoff $M_n=2$ and $t=1$,
the unique rational solution begins
\[
 B=-1+x+3x^2+16x^3+\frac{211}{2}x^4+\cdots.
\]
Indeed, $S_2(B)=1+B+B^4/2$; successively setting its coefficients of
$x^2,x^3,x^4$ to zero gives $b_2=3$, $b_3=16$, and
$-b_4+(12\cdot16+6\cdot3^2-12\cdot3+1)/2=0$.

\section{Source coverage and verification}
The following table covers all the conjectures and cross-references in these
three entries. Each COMMENT contains the same integrality conjecture, and each
cross-reference points to the other two entries.
\begin{center}\small
\begin{tabular}{@{}>{\raggedright\arraybackslash}p{.21\linewidth}>{\raggedright\arraybackslash}p{.29\linewidth}>{\raggedright\arraybackslash}p{.41\linewidth}@{}}
\toprule
Source item & Normalization and cutoff & Coverage\\
\midrule
A350410, rev.\ 5, COMMENT & $B=-1+x+\cdots$, $M_n=n$ & Proved: integrality; existence and uniqueness included.\\[3pt]
A350411, rev.\ 5, COMMENT & $B=-1+x+\cdots$, $M_n=n+1$ & Proved: integrality; existence and uniqueness included.\\[3pt]
A350412, rev.\ 9, COMMENT & $B=A-1=-1+x+\cdots$, $M_n=2n$ & Proved: integrality; existence and uniqueness included.\\[3pt]
EXAMPLE and PROGRAM sections & Finite coefficient data and recursions & Recursion~\eqref{eq:recursion} explains the formulas; coefficient data admit exact finite checks.\\[3pt]
A350412, keyword \texttt{nonn} & Nonnegative terms & Nonnegativity beyond the checked prefix is not covered.\\[3pt]
General gcd criterion and scaling & Theorem~\ref{thm:main}, Corollary~\ref{cor:scaling} & Not posed; proved for all degrees.\\[3pt]
Universal cancellation & Theorem~\ref{thm:sharp} & Not posed; exact arithmetic criterion proved.\\[3pt]
Every positive cutoff & $M_n=2$ & Not asserted in the source; fails at $b_4=211/2$.\\
\bottomrule
\end{tabular}
\end{center}
In the pinned \href{https://github.com/the-omega-institute/trureturing/blob/d75bdd6e258ca112f6f8d8ddce3c6a962cc2049f/D5/S1/Recurrence/Algebraic/HannaSquareExponentIntegrality.lean}{Lean source},
in namespace \texttt{D5.S1.Recurrence.Algebraic}, the declaration
\texttt{HannaSquareExponentIntegrality.result} proves existence,
uniqueness and integrality for exactly these three normalized rational series.
Its source predicates are \texttt{isSourceA350410} for cutoff $n$,
\texttt{isSourceA350411} for $n+1$, and \texttt{isSourceA350412} for $2n$,
with $A=B+1$ in the third case.
The same file's private \texttt{all\_for\_cutoff} theorem also covers
degree-by-degree mixtures of $n$, $n+1$ and $2n$, with linear coefficient $1$.
The general greatest-common-divisor criterion, arbitrary rational linear
coefficient, scaling and sharpness are established by the written proofs here.
The accompanying \texttt{verify.py} computes both the rational coefficient
recursion and its integer cancellation formula, compares the listed OEIS terms,
and checks the defining equations through degree $16$ using exact arithmetic.
It also matches all 173 displayed power-table coefficients, checks scaling
through degree $10$, and checks the sharpness witnesses for
$1\leq n\leq48$, $1\leq M\leq32$.
These finite checks illustrate the identities; the theorem holds in all degrees
by the induction above.

\begingroup\small
\begin{thebibliography}{9}
\bibitem{hanna410} P. D. Hanna,
\emph{A350410}, The On-Line Encyclopedia of Integer Sequences,
December 29, 2021; revision 5, December 30, 2021.
NAME, COMMENT, EXAMPLE and PROGRAM.
\url{https://oeis.org/history?seq=A350410}.
\bibitem{hanna411} P. D. Hanna,
\emph{A350411}, The On-Line Encyclopedia of Integer Sequences,
December 29, 2021; revision 5, December 30, 2021.
NAME, COMMENT, EXAMPLE and PROGRAM.
\url{https://oeis.org/history?seq=A350411}.
\bibitem{hanna412} P. D. Hanna,
\emph{A350412}, The On-Line Encyclopedia of Integer Sequences,
December 29, 2021; revision 9, December 30, 2021.
NAME, COMMENT, EXAMPLE and PROGRAM.
\url{https://oeis.org/history?seq=A350412}.
\bibitem{eger} S. Eger,
\emph{Some Elementary Congruences for the Number of Weighted Integer
Compositions}, Journal of Integer Sequences \textbf{18} (2015),
Article 15.4.1, equation (1) and Theorem 19.
\url{https://cs.uwaterloo.ca/journals/JIS/VOL18/Eger/eger11.pdf}.
\end{thebibliography}
\endgroup
\smallskip\noindent\small
\copyright\ 2026 Haobo Ma and Wenlin Zhang.
Licensed under CC BY 4.0.
\end{document}
